\documentclass[11pt,a4paper]{article}
\usepackage[utf8]{inputenc}
\usepackage[T1]{fontenc}
\usepackage{lmodern}
\usepackage[french]{babel}
\usepackage{amsmath,amssymb,mathtools}
\usepackage{icomma}
\usepackage{textcomp}
\usepackage{xcolor}
\usepackage[most]{tcolorbox}
\usepackage{enumitem}
\usepackage[a4paper,margin=2.2cm,headheight=26pt,headsep=10pt]{geometry}
\usepackage{fancyhdr}
\usepackage[hidelinks]{hyperref}
\definecolor{primary}{RGB}{222,49,99}
\definecolor{primarydark}{RGB}{150,22,55}
\tcbset{boxbase/.style={enhanced,breakable,boxrule=0.9pt,arc=2.5pt,
  left=3mm,right=3mm,top=2.3mm,bottom=2.3mm,fonttitle=\bfseries,coltitle=white}}
\newtcolorbox[auto counter]{exo}[1][]{boxbase,colback=primary!4,colframe=primary,
  title={\textsc{Exercice}~\thetcbcounter\ifx\relax#1\relax\relax\else\textmd{ --- \itshape #1}\fi}}
\newcommand{\R}{\mathbb{R}}
\newcommand{\euro}{\texteuro}
\newcommand{\ds}{\displaystyle}
\pagestyle{fancy}\fancyhf{}
\fancyhead[L]{\small\textcolor{primarydark}{\textbf{Thème 4} — Systèmes d'équations}}
\fancyhead[R]{\small\textcolor{primarydark}{\textsc{Facile}}}
\fancyfoot[C]{\small\thepage}
\renewcommand{\headrulewidth}{0.8pt}
\begin{document}

\begin{center}
{\LARGE\bfseries\color{primarydark} Niveau Facile — Exercices}\\[2pt]
{\color{primary}\rule{0.45\textwidth}{1.2pt}}\\[3pt]
{\small\itshape Une seule mécanique à la fois. Aucune calculatrice.}
\end{center}
\vspace{1mm}

\begin{exo}[résoudre par substitution]
Une inconnue est déjà isolée : remplace-la dans l'autre équation.
\[ \text{a) }\begin{cases} y=x+1\\ 2x+y=7\end{cases}\qquad
   \text{b) }\begin{cases} x=2y\\ x+3y=10\end{cases}\qquad
   \text{c) }\begin{cases} y=3\\ x+y=5\end{cases} \]
\end{exo}

\begin{exo}[résoudre par combinaison]
Additionne ou soustrais les deux équations pour éliminer une inconnue :
\[ \text{a) }\begin{cases} x+y=6\\ x-y=2\end{cases}\qquad
   \text{b) }\begin{cases} 2x+y=7\\ x-y=2\end{cases}\qquad
   \text{c) }\begin{cases} 3x+2y=12\\ 3x-2y=0\end{cases} \]
\end{exo}

\begin{exo}[vérifier une solution]
Le couple $(x\,;y)=(2\,;-1)$ est-il solution de chacun de ces systèmes~? Justifie en remplaçant.
\[ \text{a) }\begin{cases} x+y=1\\ x-y=3\end{cases}\qquad
   \text{b) }\begin{cases} 2x+y=3\\ x+2y=1\end{cases}\qquad
   \text{c) }\begin{cases} x+y=1\\ 3x-y=5\end{cases} \]
\end{exo}

\begin{exo}[traduire un énoncé en système]
Écris (sans résoudre) le système d'équations correspondant :
\begin{enumerate}[label=\alph*),leftmargin=2em,topsep=1pt,itemsep=1pt]
  \item La somme de deux nombres $x$ et $y$ vaut $20$ ; leur différence vaut $4$.
  \item Un stylo ($x$) et un cahier ($y$) coûtent $5$~\euro{} ; deux stylos et un cahier coûtent $7$~\euro.
  \item $x$ est le triple de $y$ ; leur somme vaut $16$.
\end{enumerate}
\end{exo}

\begin{exo}[résoudre un système $2\times2$ (méthode libre)]
\[ \text{a) }\begin{cases} x+2y=5\\ 3x+y=20\end{cases}\qquad\qquad
   \text{b) }\begin{cases} 2x+3y=8\\ x-y=-1\end{cases} \]
\end{exo}

\end{document}
